% Lösungen Einheit 4 – Ausklammern (zusammenbau v1.0)
\kbabschnitt{Ausklammern}

\erg{20}{$5x \cdot (x + 2y - 3)$; Probe: $5x \cdot x + 5x \cdot 2y - 5x \cdot 3 = 5x^2 + 10xy - 15x$ (richtig $\rightarrow$ Nr.~22–25 überspringen)}
\erg{21}{a) die Zahl 4 \quad b) die Variable $y$ \quad c) die Zahl 2 und die Variable $a$, also $2a$ \quad d) die Zahl 3}
\erg{22}{a) In die Lücken: $2x$ und $5$, also $8x + 20 = 4 \cdot 2x + 4 \cdot 5$ \quad b) In die Lücken: $3a$ und $7$, also $6a + 14 = 2 \cdot 3a + 2 \cdot 7$ \quad c) In die Lücken: $2y$ und $3$, also $10y + 15 = 5 \cdot 2y + 5 \cdot 3$ \quad d) In die Lücken: $3b$ und $2$, also $9b + 6 = 3 \cdot 3b + 3 \cdot 2$}
\erg{23}{a) In die Klammer: $3x$ und $7$, also $9x + 21 = 3 \cdot (3x + 7)$ \quad b) In die Klammer: $2y$ und $5$, also $4y + 10 = 2 \cdot (2y + 5)$ \quad c) In die Klammer: $2a$ und $7$, also $10a + 35 = 5 \cdot (2a + 7)$ \quad d) In die Klammer: $2b$ und $3$, also $14b + 21 = 7 \cdot (2b + 3)$}
\erg{24}{a) Faktor finden: $5$; ausklammern: $5x + 10 = 5 \cdot (x + 2)$ \quad b) Faktor finden: $5$; ausklammern: $5x + 25 = 5 \cdot (x + 5)$ \quad c) Faktor finden: $5$; ausklammern: $5x + 30 = 5 \cdot (x + 6)$ \quad d) Faktor finden: $5$; ausklammern: $5x + 35 = 5 \cdot (x + 7)$ \quad e) $x \cdot (x + 5)$ \quad f) $3x \cdot (2x + 5)$ \quad g) $6 \cdot (x + 1)$ \quad h) $3 \cdot (x^2 + 2x + 5)$ \quad i) $6 \cdot (x + 5)$; Probe: $6 \cdot x + 6 \cdot 5 = 6x + 30$ \quad j) Die 20 wurde nicht durch 5 geteilt, die Probe ergibt $5x + 100$. Richtig: $5x + 20 = 5 \cdot (x + 4)$.}
\erg{25}{$3a \cdot (a + 2b + 3)$; Probe: $3a \cdot a + 3a \cdot 2b + 3a \cdot 3 = 3a^2 + 6ab + 9a$}
\erg{26}{a) Das Minus ging verloren, die Probe ergibt $20x + 8$. Richtig: Faktor finden: $4$; ausklammern: $20x - 8 = 4 \cdot (5x - 2)$. \quad b) Distributivgesetz: ausmultipliziert ist $6 \cdot (a + 7) = 6a + 42$, beide Terme sind gleichwertig; Probe mit $a = 2$: $6 \cdot 2 + 42 = 54$ und $6 \cdot 9 = 54$. \quad c) Zwei Rechtecke nebeneinander, beide 6 cm hoch, eines $b$ cm breit, eines 3 cm breit; zusammen $6b + 18$ cm². \quad d) Ja; Term mit Klammer: $15 \cdot (28 + 17)$; rechnen: $15 \cdot 45 = 675$ €, das sind 25 € weniger als 700 €.}
